\documentclass[11pt,a4paper]{article}

\usepackage[T1]{fontenc}
\usepackage[utf8]{inputenc}
\usepackage[brazilian]{babel}
\usepackage{amsmath,amssymb,mathtools}
\usepackage{hyperref}
\usepackage[most]{tcolorbox}
\usepackage{geometry}
\usepackage{enumitem}
\usepackage{xcolor}
\geometry{margin=2.4cm}
\hypersetup{colorlinks=true,linkcolor=blue!60!black,urlcolor=blue!60!black}
\setlist[itemize]{topsep=4pt,itemsep=2pt}
\setlength{\parskip}{0.5em}
\setlength{\parindent}{0pt}
\newcommand{\dd}{\mathrm{d}}
\newcommand{\E}{\mathrm{e}}

\newtcolorbox{resultado}{
  colback=blue!4!white,
  colframe=blue!55!black,
  title=Resultado final,
  breakable
}

\newtcolorbox{intuicao}{
  colback=green!4!white,
  colframe=green!45!black,
  title=Intui\c{c}\~ao econ\^omica,
  breakable
}

\newtcolorbox{convencao}{
  colback=gray!8!white,
  colframe=gray!55!black,
  title=Conven\c{c}\~oes desta nota,
  breakable
}

\title{Cap\'itulo 2 --- Deriva\c{c}\~oes do RCK e do Diamond}
\author{romer-study}
\date{\today}

\begin{document}

\maketitle
\tableofcontents
\bigskip

\section{Objetivo da nota}

Esta nota cobre o Cap\'itulo 2 do Romer em duas partes:
\begin{itemize}
  \item o modelo de Ramsey--Cass--Koopmans (RCK), em tempo cont\'inuo;
  \item o modelo de Diamond, com gera\c{c}\~oes sobrepostas (OLG), em tempo discreto.
\end{itemize}

O objetivo \'e mostrar a deriva\c{c}\~ao formal das equa\c{c}\~oes, ligar a
nota\c{c}\~ao do livro \`a implementa\c{c}\~ao do reposit\'orio e deixar claro
onde entram as condi\c{c}\~oes de steady state, as is\'oclinas e a discuss\~ao
de inefici\^encia din\^amica.

\begin{convencao}
\begin{itemize}
  \item Na parte do RCK, usamos \(k\) e \(c\) como vari\'aveis por trabalhador efetivo.
  \item O c\'odigo do projeto inclui \(\delta>0\) e um gasto do governo \(G\) por trabalhador efetivo.
  \item Na parte do Diamond, usamos utilidade logar\'itmica com peso intertemporal \(\beta\).
\end{itemize}
\end{convencao}

\section{Parte A --- Ramsey--Cass--Koopmans}

\subsection{Ambiente econ\^omico}

A tecnologia agregada \'e
\[
Y(t)=F(K(t),A(t)L(t))=A(t)L(t)f(k(t)),
\qquad
k(t)\equiv \frac{K(t)}{A(t)L(t)}.
\]
As taxas de crescimento ex\'ogeno s\~ao
\[
\frac{\dot L}{L}=n,
\qquad
\frac{\dot A}{A}=g.
\]

A restri\c{c}\~ao agregada de recursos \'e
\[
\dot K = Y - C - \mathcal G - \delta K,
\]
onde \(\mathcal G\) denota o gasto agregado do governo.

\subsection{Normaliza\c{c}\~ao da restri\c{c}\~ao}

Defina
\[
c\equiv \frac{C}{AL},
\qquad
G\equiv \frac{\mathcal G}{AL},
\qquad
k\equiv \frac{K}{AL}.
\]
Derivando \(k=K/(AL)\),
\[
\dot k
=
\frac{\dot K}{AL}
-
\frac{K}{(AL)^2}\frac{\dd(AL)}{\dd t}.
\]
Como \(\dd(AL)/\dd t = AL(n+g)\), segue que
\[
\dot k=\frac{\dot K}{AL}-(n+g)k.
\]

Agora dividimos a restri\c{c}\~ao agregada por \(AL\):
\[
\frac{\dot K}{AL}=f(k)-c-G-\delta k.
\]
Substituindo,
\[
\dot k=f(k)-c-G-\delta k-(n+g)k.
\]
Portanto,
\[
\dot k=f(k)-c-G-(n+g+\delta)k.
\]

\begin{resultado}
No sistema usado no reposit\'orio,
\[
\dot k=f(k)-c-G-(n+g+\delta)k.
\]
\end{resultado}

\subsection{Normaliza\c{c}\~ao da utilidade}

A fam\'ilia representativa maximiza
\[
U=\int_0^\infty \E^{-\rho t}L(t)\frac{\widetilde c(t)^{1-\theta}}{1-\theta}\,\dd t,
\]
onde \(\widetilde c(t)\equiv C(t)/L(t)\) \'e o consumo por trabalhador.

Como \(\widetilde c = A c\), temos
\[
U
=
\int_0^\infty
\E^{-\rho t}L(t)\frac{(A(t)c(t))^{1-\theta}}{1-\theta}\,\dd t.
\]
Usando
\[
L(t)=L(0)\E^{nt},
\qquad
A(t)=A(0)\E^{gt},
\]
segue que
\[
U
=
\frac{L(0)A(0)^{1-\theta}}{1-\theta}
\int_0^\infty
\E^{-\rho t}\E^{nt}\E^{(1-\theta)gt}c(t)^{1-\theta}\,\dd t.
\]
Reagrupando os expoentes,
\[
U
=
\text{constante}
\cdot
\int_0^\infty \E^{-\eta t}\frac{c(t)^{1-\theta}}{1-\theta}\,\dd t,
\]
com
\[
\eta \equiv \rho - n - (1-\theta)g.
\]

\begin{resultado}
Depois da normaliza\c{c}\~ao por trabalho efetivo, o problema passa a ser
\[
\max \int_0^\infty \E^{-\eta t}\frac{c(t)^{1-\theta}}{1-\theta}\,\dd t,
\qquad
\eta=\rho-n-(1-\theta)g.
\]
\end{resultado}

\begin{intuicao}
O ajuste da taxa de desconto aparece porque o problema original est\'a em n\'ivel
agregado e precisa ser reescrito em unidades por trabalhador efetivo.
\end{intuicao}

\subsection{Hamiltoniano e condi\c{c}\~oes de primeira ordem}

O Hamiltoniano em valor corrente pode ser escrito como
\[
\mathcal H
=
\frac{c^{1-\theta}}{1-\theta}
\lambda \left[f(k)-c-G-(n+g+\delta)k\right].
\]

As condi\c{c}\~oes de primeira ordem s\~ao:

\paragraph{FOC em rela\c{c}\~ao a \(c\).}
\[
\frac{\partial \mathcal H}{\partial c}
=
c^{-\theta}-\lambda=0
\qquad \Longrightarrow \qquad
\lambda=c^{-\theta}.
\]

\paragraph{Equa\c{c}\~ao do coestado.}
\[
\dot\lambda
=
\lambda\left[\eta-(f'(k)-(n+g+\delta))\right].
\]

Substituindo \(\eta=\rho-n-(1-\theta)g\),
\[
\dot\lambda
=
\lambda\left[\rho+\delta+\theta g-f'(k)\right].
\]

\subsection{Deriva\c{c}\~ao da equa\c{c}\~ao de Euler}

Da FOC, \(\lambda=c^{-\theta}\). Derivando no tempo,
\[
\dot\lambda
=
-\theta c^{-\theta-1}\dot c.
\]
Dividindo ambos os lados por \(\lambda=c^{-\theta}\),
\[
\frac{\dot\lambda}{\lambda}
=
-\theta \frac{\dot c}{c}.
\]

Da equa\c{c}\~ao do coestado,
\[
\frac{\dot\lambda}{\lambda}
=
\rho+\delta+\theta g-f'(k).
\]
Igualando as duas express\~oes,
\[
-\theta \frac{\dot c}{c}
=
\rho+\delta+\theta g-f'(k).
\]
Multiplicando por \(-1\),
\[
\theta \frac{\dot c}{c}
=
f'(k)-\delta-\rho-\theta g.
\]
Logo,
\[
\frac{\dot c}{c}
=
\frac{f'(k)-\delta-\rho-\theta g}{\theta}.
\]

\begin{resultado}
A equa\c{c}\~ao de Euler do RCK, na conven\c{c}\~ao do projeto, \'e
\[
\frac{\dot c}{c}
=
\frac{f'(k)-\delta-\rho-\theta g}{\theta}.
\]
\end{resultado}

\subsection{Sistema din\^amico bidimensional}

Juntando a restri\c{c}\~ao din\^amica e a equa\c{c}\~ao de Euler, obtemos
\[
\dot k=f(k)-c-G-(n+g+\delta)k,
\]
\[
\dot c=
c\cdot
\frac{f'(k)-\delta-\rho-\theta g}{\theta}.
\]
Esse \'e um sistema bidimensional em \((k,c)\) porque tanto o estoque de
capital quanto o consumo s\~ao vari\'aveis end\'ogenas de estado/controle na
trajet\'oria de equil\'ibrio.

\subsection{Is\'oclinas}

\paragraph{\(\dot c=0\).}
Como \(c>0\), temos
\[
\dot c=0
\qquad \Longleftrightarrow \qquad
f'(k)=\delta+\rho+\theta g.
\]
Portanto, a is\'oclina \(\dot c=0\) \'e uma linha vertical em
\[
k=k^{\ast\ast},
\]
onde \(k^{\ast\ast}\) satisfaz essa igualdade.

\paragraph{\(\dot k=0\).}
\[
\dot k=0
\qquad \Longleftrightarrow \qquad
c=f(k)-G-(n+g+\delta)k.
\]
Essa \'e uma curva decrescente em \(c\) para cada \(k\), e no c\'odigo aparece
como a fun\c{c}\~ao \texttt{k\_locus(...)}.

\subsection{Steady state}

No steady state, \(\dot k=0\) e \(\dot c=0\). Logo,
\[
f'(k^{\ast\ast})=\rho+\theta g+\delta,
\]
\[
c^{\ast\ast}=f(k^{\ast\ast})-G-(n+g+\delta)k^{\ast\ast}.
\]

\begin{resultado}
A condi\c{c}\~ao de steady state do RCK \'e
\[
f'(k^{\ast\ast})=\rho+\theta g+\delta.
\]
\end{resultado}

No caso Cobb--Douglas, \(f(k)=k^\alpha\), logo
\[
f'(k)=\alpha k^{\alpha-1}.
\]
Portanto,
\[
\alpha (k^{\ast\ast})^{\alpha-1}=\rho+\theta g+\delta
\]
e
\[
k^{\ast\ast}
=
\left(\frac{\alpha}{\rho+\theta g+\delta}\right)^{\frac{1}{1-\alpha}}.
\]

\subsection{Trajet\'oria de sela}

O steady state do RCK n\~ao \'e um ponto atrator em duas dimens\~oes como no
Solow unidimensional. A estabilidade ocorre ao longo de uma \emph{trajet\'oria
de sela}: apenas um valor inicial de \(c(0)\) para cada \(k(0)\) coloca a
economia no caminho convergente.

Se \(c(0)\) for alto demais, o capital colapsa. Se \(c(0)\) for baixo demais,
o capital tende a divergir. O algoritmo de shooting em
\texttt{ch02\_rck\_diamond/ch02\_rck.py} procura exatamente esse \(c(0)\)
compat\'ivel com a trajet\'oria est\'avel.

\begin{intuicao}
O RCK endogeneiza a poupan\c{c}a. Por isso, o consumo inicial n\~ao pode ser
escolhido arbitrariamente: ele precisa ser consistente com a trajet\'oria que
satisfaz a condi\c{c}\~ao de transversalidade.
\end{intuicao}

\subsection{Condi\c{c}\~ao de transversalidade}

A condi\c{c}\~ao de transversalidade elimina trajet\'orias em que a fam\'ilia
acumula ativos sem motivo econ\^omico. Em linguagem econ\^omica, ela impede
que valor seja deixado infinitamente ``sobre a mesa'' no limite.

Na formula\c{c}\~ao normalizada, a ideia \'e que o valor sombra do capital,
ponderado adequadamente, deve convergir a zero no infinito. Essa condi\c{c}\~ao
\'e o crit\'erio conceitual por tr\'as da sele\c{c}\~ao da saddle path.

\subsection{Do livro para o c\'odigo}

As correspond\^encias principais com \texttt{ch02\_rck\_diamond/ch02\_rck.py}
s\~ao:
\begin{itemize}
  \item \(f(k)\) \(\leftrightarrow\) \texttt{f(...)};
  \item \(f'(k)\) \(\leftrightarrow\) \texttt{f\_prime(...)};
  \item \(\dot k\) \(\leftrightarrow\) \texttt{k\_dot(...)};
  \item \(\dot c\) \(\leftrightarrow\) \texttt{c\_dot(...)};
  \item \(\dot k=0\) \(\leftrightarrow\) \texttt{k\_locus(...)};
  \item \((k^{\ast\ast},c^{\ast\ast})\) \(\leftrightarrow\) \texttt{steady\_state(...)};
  \item a trajet\'oria est\'avel \(\leftrightarrow\) \texttt{find\_saddle\_path(...)}.
\end{itemize}

\section{Parte B --- O modelo de Diamond}

\subsection{Estrutura b\'asica}

No modelo de Diamond, os agentes vivem por dois per\'iodos:
\begin{itemize}
  \item quando jovens, trabalham, recebem sal\'ario e poupam;
  \item quando velhos, n\~ao trabalham e consomem o retorno da poupan\c{c}a.
\end{itemize}

Seja \(c_{1,t}\) o consumo quando jovem e \(c_{2,t+1}\) o consumo quando velho.
A utilidade da gera\c{c}\~ao nascida em \(t\) \'e
\[
u_t=\ln(c_{1,t})+\beta \ln(c_{2,t+1}).
\]

\subsection{Restri\c{c}\~oes or\c{c}ament\'arias}

Se \(w_t\) \'e o sal\'ario recebido quando jovem e \(s_t\) a poupan\c{c}a,
ent\~ao
\[
c_{1,t}+s_t=w_t.
\]
Quando velho,
\[
c_{2,t+1}=(1+r_{t+1})s_t.
\]
Substituindo na utilidade,
\[
u_t
=
\ln(w_t-s_t)+\beta \ln\!\bigl((1+r_{t+1})s_t\bigr).
\]

\subsection{Regra de poupan\c{c}a do jovem}

O problema da gera\c{c}\~ao jovem \'e
\[
\max_{s_t}
\ln(w_t-s_t)+\beta \ln\!\bigl((1+r_{t+1})s_t\bigr).
\]
A condi\c{c}\~ao de primeira ordem \'e
\[
-\frac{1}{w_t-s_t}+\beta \frac{1}{s_t}=0.
\]
Reorganizando,
\[
\beta(w_t-s_t)=s_t.
\]
Logo,
\[
\beta w_t=(1+\beta)s_t
\qquad \Longrightarrow \qquad
s_t=\frac{\beta}{1+\beta}w_t.
\]

\begin{resultado}
Com utilidade logar\'itmica, a poupan\c{c}a do jovem \'e uma fra\c{c}\~ao fixa do sal\'ario:
\[
s_t=\frac{\beta}{1+\beta}w_t.
\]
\end{resultado}

\subsection{Do sal\'ario ao mapa de capital}

Suponha produ\c{c}\~ao Cobb--Douglas por trabalhador,
\[
y_t=k_t^\alpha,
\]
e deprecia\c{c}\~ao completa, \(\delta=1\), como no caso padr\~ao mais simples do OLG.
Ent\~ao o sal\'ario \'e
\[
w_t=f(k_t)-k_t f'(k_t).
\]
Para Cobb--Douglas,
\[
f'(k_t)=\alpha k_t^{\alpha-1},
\]
logo
\[
w_t
=
k_t^\alpha-\alpha k_t^\alpha
=
(1-\alpha)k_t^\alpha.
\]

Se a popula\c{c}\~ao cresce \`a taxa \(n\), ent\~ao o capital do per\'iodo
seguinte por jovem \'e a poupan\c{c}a agregada dividida pela nova coorte:
\[
k_{t+1}=\frac{s_t}{1+n}.
\]
Substituindo a regra de poupan\c{c}a,
\[
k_{t+1}
=
\frac{1}{1+n}\frac{\beta}{1+\beta}w_t
=
\frac{1}{1+n}\frac{\beta}{1+\beta}(1-\alpha)k_t^\alpha.
\]

\begin{resultado}
O mapa din\^amico do Diamond com Cobb--Douglas e utilidade log \'e
\[
k_{t+1}
=
\frac{\beta(1-\alpha)}{(1+\beta)(1+n)}k_t^\alpha.
\]
\end{resultado}

\subsection{Steady state do mapa discreto}

Um steady state satisfaz
\[
k^\ast = \frac{\beta(1-\alpha)}{(1+\beta)(1+n)}(k^\ast)^\alpha.
\]
Se \(k^\ast>0\), dividimos por \((k^\ast)^\alpha\):
\[
(k^\ast)^{1-\alpha}
=
\frac{\beta(1-\alpha)}{(1+\beta)(1+n)}.
\]
Logo,
\[
k^\ast
=
\left[
\frac{\beta(1-\alpha)}{(1+\beta)(1+n)}
\right]^{\frac{1}{1-\alpha}}.
\]

\subsection{Inefici\^encia din\^amica e compara\c{c}\~ao com a Regra de Ouro}

No OLG, pode acontecer de o steady state competitivo gerar capital demais. A
refer\^encia \'e o capital da Regra de Ouro, definido por
\[
f'(k_{\text{gold}})=n+\delta.
\]
No caso simples com \(\delta=1\) na formula\c{c}\~ao discreta padr\~ao, a regra
equivalente \'e comparar o produto marginal l\'iquido do capital com a taxa de
crescimento da economia. Intuitivamente:
\begin{itemize}
  \item se \(k^\ast < k_{\text{gold}}\), a economia poupa pouco;
  \item se \(k^\ast > k_{\text{gold}}\), h\'a sobreacumula\c{c}\~ao e inefici\^encia din\^amica.
\end{itemize}

\begin{intuicao}
O ponto central do Diamond \'e que equil\'ibrio competitivo n\~ao garante
efici\^encia intertemporal. Diferentemente do RCK, a estrutura demogr\'afica
pode levar a um estoque de capital excessivo.
\end{intuicao}

\section{Fechamento}

O Cap\'itulo 2 do Romer amplia o Solow em duas dire\c{c}\~oes:
\begin{itemize}
  \item no RCK, a poupan\c{c}a \'e end\'ogena e a din\^amica passa a ser um sistema em \((k,c)\);
  \item no Diamond, a demografia e a sobreposi\c{c}\~ao de gera\c{c}\~oes tornam poss\'ivel a inefici\^encia din\^amica.
\end{itemize}

No reposit\'orio, o RCK j\'a aparece implementado em
\texttt{ch02\_rck\_diamond/ch02\_rck.py}. O Diamond, por enquanto, fica como
ponte te\'orica para o pr\'oximo m\'odulo.

\end{document}
